\begin{abstract}
Rapid growth of behind-the-meter photovoltaic generation reshapes the load profile seen by transmission system operators, making day-ahead load forecasting a problem of persistent concept drift. Although open data now allow research models to be compared with the forecasts actually used in operation, such comparisons are rare. This paper presents an open and fully reproducible benchmark for day-ahead hourly load forecasting of the four subsystems of the Brazilian interconnected system, where behind-the-meter generation grew from about 5\% to one third of the midday load in five years. Using only public operator data (2012--2026), forecasts issued at noon of the previous day, and a test period of almost three years, we compare seasonal naive benchmarks, hour-specific linear autoregressive models, gradient boosting, a neural network, and their ensemble; study calibration windows and refitting frequencies; compare quantile regression with adaptive conformal prediction intervals; and evaluate the models against the operator's official day-ahead forecast. Without any weather input, the ensemble achieved a mean absolute percentage error of 2.44\%, versus 2.90\% for the official forecast, with significant improvements in three of the four subsystems according to Diebold--Mariano tests. The official forecast showed a systematic midday under-forecast that deepened from 2024 to 2026 as behind-the-meter generation grew, whereas the open models remained nearly unbiased. Level normalisation was essential for tree-based models, the calibration window mattered more than the refitting frequency, and adaptive conformal inference produced calibrated intervals where quantile gradient boosting was overconfident. Monitoring hour-specific bias and maintaining data-driven benchmarks can help operators cope with distributed solar growth.
\end{abstract}

\begin{keywords}
Load forecasting, behind-the-meter photovoltaics, concept drift, conformal prediction, gradient boosting, power system operation, Brazil.
\end{keywords}
